\scanpage{69}
% archive sheet number: 68
% This scan is handwritten throughout; the faint mirrored typescript is show-through from the reverse, not a second front-side text.
$\overline{\text{\unsure{un|Un}}}$ énoncé pour $X=\operatorname{Spec}\mathbb{Z}$, et dans le cas géométrique, pour $X=\operatorname{Spec}\mathbb{Z}/p\mathbb{Z}$. Dans ce dernier cas, on doit vérifier simplement la dualité entre
\[
H^i(\mathbb{F}_p,M)\quad\text{et}\quad H^{1-i}(\mathbb{F}_p,M^*)
\]
\[
(M^*=\unsure{\operatorname{Hom}|\mathscr{H}\mathrm{om}}(M,\unsure{\mathbb{Q}(0)|\mathbb{Q}\omega)}))
\]
à l'aide des ``\uline{cup-produits}'' et de l'isom.
\margin{\unsure{à développer un peu pour les cup-produits à cohomologie ``absolue''} $-\ -\ -$}
\[
H^1(\mathbb{F}_p,\mathbb{Q}(0))\simeq\mathbb{Q}
\]
et il suffit de faire la \unsure{vérification} $l$-adiquement, en appliquant le foncteur $T_\ell$. Cela revient : \unsure{la|le} \unsure{vérification suivante} [pour un \unsure{vectoriel|vecteur} de dim finie sur $\mathbb{Q}_\ell$ sur lequel $\pi=\mathbb{Z}$ opère]
\[
H^i(\pi,V)\quad\text{et}\quad H^{1-i}(\pi,V^*)\qquad\text{en dualité}
\]
\[
(V^*=\operatorname{Hom}_{\mathbb{Q}_\ell}(V,\mathbb{Q}_\ell)\ :\ \text{c'est une}
\]
trivialité.
